% =============================================================================
% 01_pointer_memory_model.tex — Slides S01-005 to S01-026
% Block 1: The Object Pointer Mental Model & Memory Semantics
% =============================================================================

% -----------------------------------------------------------------------------
% Slide ID: S01-005
% Narrative Role: Block Transition / Tension
% Cognitive Objective: Establish the local mystery of Block 1.
% Plan Source: slide_plan.md / S01-005
% -----------------------------------------------------------------------------
\TopicDivider{01}{The Object Pointer Mental Model}{Stack pointers, heap PyObjects, mutability traps, and copy semantics.}

% -----------------------------------------------------------------------------
% Slide ID: S01-006
% Narrative Role: Prediction Prompt
% Cognitive Objective: Commit learners to a prediction regarding mutable default arguments.
% Plan Source: slide_plan.md / S01-006
% -----------------------------------------------------------------------------
\begin{frame}[fragile]{Prediction Task 1: The Multi-Tenant Event Logger}
  \begin{columns}[T,onlytextwidth]
    \begin{column}{0.52\textwidth}
      \begin{CodeBlock}{Python}
def append_event(event_id, log=[]):
    log.append(event_id)
    return log

# User A records an event
print("Call 1:", append_event(101))

# User B records an event
print("Call 2:", append_event(202))
      \end{CodeBlock}
    \end{column}
    \begin{column}{0.45\textwidth}
      \begin{QuestionBox}{What will Call 2 output?}
        \begin{itemize}\setlength{\itemsep}{0.12cm}
          \item[\textbf{A)}] \InlineCode{Call 1: [101] \textbar\ Call 2: [202]}
          \item[\textbf{B)}] \InlineCode{Call 1: [101] \textbar\ Call 2: [101, 202]}
          \item[\textbf{C)}] \InlineCode{Call 1: [101] \textbar\ Call 2: None}
          \item[\textbf{D)}] \InlineCode{TypeError: unhashable default}
        \end{itemize}
      \end{QuestionBox}
      \vspace{0.10cm}
      {\scriptsize\color{UpGradSlate}Commit your prediction in the live poll before proceeding.}
    \end{column}
  \end{columns}
\end{frame}

% -----------------------------------------------------------------------------
% Slide ID: S01-007
% Narrative Role: Reveal / Dissonance
% Cognitive Objective: Expose the counter-intuitive reality that User B received User A's data.
% Plan Source: slide_plan.md / S01-007
% -----------------------------------------------------------------------------
\begin{frame}[fragile]{Reveal 1: The Audit Log Contamination Bug}
  \begin{columns}[T,onlytextwidth]
    \begin{column}{0.48\textwidth}
      \begin{CodeBlock}{Python}
def append_event(event_id, log=[]):
    log.append(event_id)
    return log

append_event(101)  # User A
append_event(202)  # User B
      \end{CodeBlock}
    \end{column}
    \begin{column}{0.48\textwidth}
      \begin{TerminalWindow}[PYTHON 3.12 INTERPRETER]
>>> Call 1: [101] \\
>>> Call 2: [101, 202]  \# <--- CRITICAL STATE LEAK!
      \end{TerminalWindow}
      \vspace{0.08cm}
      \begin{WarningBox}{Cross-Tenant Data Leak}
        \scriptsize User B received User A's private audit telemetry. Why did Call 2 share User A's list rather than creating a fresh one?
      \end{WarningBox}
    \end{column}
  \end{columns}
  \vspace{0.08cm}
  \TakeawayBanner{Option B is correct. Python default arguments are evaluated once at definition time!}
\end{frame}

% -----------------------------------------------------------------------------
% Slide ID: S01-008
% Narrative Role: Misconception Breakdown
% Cognitive Objective: Break the "variable as a value box" mental model.
% Plan Source: slide_plan.md / S01-008
% -----------------------------------------------------------------------------
\begin{frame}{Naive Mental Model vs Physical Reality}
  \CompareGrid{Naive Model: The Variable as a Box}{
    \begin{itemize}\setlength{\itemsep}{0.08cm}
      \item $x$ is a container labeled \InlineCode{"x"}.
      \item Assigning $x = 42$ puts 42 inside the box.
      \item Assigning $y = x$ creates a second box with a copy of 42.
      \item Modifying $y$ never touches $x$.
    \end{itemize}
  }{Physical Reality: Names Point to Heap PyObjects}{
    \begin{itemize}\setlength{\itemsep}{0.08cm}
      \item $x$ is an 8-byte pointer in the stack frame.
      \item 42 lives in a 28-byte \InlineCode{PyObject} struct in heap RAM.
      \item Assigning $y = x$ copies the 64-bit pointer address.
      \item Both names reference the \textbf{exact same heap struct}!
    \end{itemize}
  }
  \vspace{0.18cm}
  \TakeawayBanner{In Python, variables store memory addresses, never raw data values.}
\end{frame}

% -----------------------------------------------------------------------------
% Slide ID: S01-009
% Narrative Role: Mechanistic Explanation
% Cognitive Objective: Understand the internal C structure of all Python objects.
% Plan Source: slide_plan.md / S01-009
% -----------------------------------------------------------------------------
\begin{frame}{Inside the Heap: Anatomy of a PyObject}
  \begin{columns}[c,onlytextwidth]
    \begin{column}{0.54\textwidth}
      \DrawPyObjectStruct{0x100A}{1}{<class 'int'>}{42}
    \end{column}
    \begin{column}{0.43\textwidth}
      \begin{ConceptBox}{The 28-Byte Overhead}
        \scriptsize
        \begin{itemize}\setlength{\itemsep}{0.05cm}
          \item \textbf{ob\_refcnt} (8B): Tracks active pointers.
          \item \textbf{ob\_type*} (8B): Pointer to type descriptor.
          \item \textbf{ob\_ival} (8-12B): Raw numerical payload.
        \end{itemize}
      \end{ConceptBox}
      \vspace{0.06cm}
      \StatCard{UpGradRed}{28 Bytes}{Minimum Int Size}{3.5x overhead vs 8-byte C int64}
    \end{column}
  \end{columns}
\end{frame}

% -----------------------------------------------------------------------------
% Slide ID: S01-010
% Narrative Role: Flipbook Sequence (Step 1/2)
% Cognitive Objective: Visually observe x = 42 in heap memory.
% Plan Source: slide_plan.md / S01-010
% -----------------------------------------------------------------------------
% -----------------------------------------------------------------------------
% Slide ID: S01-010
% Narrative Role: Flipbook Sequence (Step 1/2)
% Cognitive Objective: Visually observe x = 42 in heap memory.
% Plan Source: slide_plan.md / S01-010
% -----------------------------------------------------------------------------
\begin{frame}[fragile]{Variable Binding: Step 1 (\texttt{x = 42})}
  \begin{columns}[T,onlytextwidth]
    \begin{column}{0.35\textwidth}
      \begin{HighlightedCodeBlock}{Python}{1}
x = 42
      \end{HighlightedCodeBlock}
      \vspace{0.12cm}
      \begin{ConceptBox}{Execution Step 1}
        \scriptsize CPython allocates a \InlineCode{PyLongObject} at address \InlineCode{0x100A} with \InlineCode{ob\_refcnt = 1}. The stack name \InlineCode{x} stores \InlineCode{0x100A}.
      \end{ConceptBox}
    \end{column}
    \begin{column}{0.62\textwidth}
      \centering
      \DrawPointerRebinding{1}
    \end{column}
  \end{columns}
  \vspace{0.15cm}
  \TakeawayBanner{Binding assigns a memory pointer to a stack name.}
\end{frame}

% -----------------------------------------------------------------------------
% Slide ID: S01-011
% Narrative Role: Flipbook Sequence (Step 2/2)
% Cognitive Objective: Visually observe x = "hello" rebind the pointer without altering previous object.
% Plan Source: slide_plan.md / S01-011
% -----------------------------------------------------------------------------
\begin{frame}[fragile]{Variable Rebinding: Step 2 (\texttt{x = "hello"})}
  \begin{columns}[T,onlytextwidth]
    \begin{column}{0.35\textwidth}
      \begin{HighlightedCodeBlock}{Python}{2}
x = 42
x = "hello"  # Rebind!
      \end{HighlightedCodeBlock}
      \vspace{0.12cm}
      \begin{ConceptBox}{Execution Step 2}
        \scriptsize The name \InlineCode{x} is repointed to \InlineCode{0x200B}. The integer at \InlineCode{0x100A} sees \InlineCode{ob\_refcnt = 0} and is queued for garbage collection.
      \end{ConceptBox}
    \end{column}
    \begin{column}{0.62\textwidth}
      \centering
      \DrawPointerRebinding{2}
    \end{column}
  \end{columns}
  \vspace{0.15cm}
  \TakeawayBanner{Variables do not mutate types; variables repoint to different typed objects.}
\end{frame}

% -----------------------------------------------------------------------------
% Slide ID: S01-012
% Narrative Role: Formalization / Concept
% Cognitive Objective: Differentiate between content comparison (__eq__) and memory address comparison (id()).
% Plan Source: slide_plan.md / S01-012
% -----------------------------------------------------------------------------
\begin{frame}[fragile]{Value Equality (\texttt{==}) vs Reference Identity (\texttt{is})}
  \begin{columns}[T,onlytextwidth]
    \begin{column}{0.48\textwidth}
      \begin{ConceptBox}{Value Equality: \InlineCode{a == b}}
        \begin{itemize}
          \item Calls \InlineCode{a.\_\_eq\_\_(b)}.
          \item Compares internal structural contents.
          \item Returns \InlineCode{True} if values are identical.
        \end{itemize}
      \end{ConceptBox}
    \end{column}
    \begin{column}{0.48\textwidth}
      \begin{ConceptBox}{Reference Identity: \InlineCode{a is b}}
        \begin{itemize}
          \item Evaluates \InlineCode{id(a) == id(b)}.
          \item Compares 64-bit physical heap addresses.
          \item Returns \InlineCode{True} ONLY if both names share one object.
        \end{itemize}
      \end{ConceptBox}
    \end{column}
  \end{columns}
  \vspace{0.15cm}
  \begin{HighlightedCodeBlock}{Python}{3,4}
list1 = [1, 2, 3]
list2 = [1, 2, 3]
print(list1 == list2)  # True  (Contents match!)
print(list1 is list2)  # False (Two distinct physical heap allocations!)
  \end{HighlightedCodeBlock}
\end{frame}

% -----------------------------------------------------------------------------
% Slide ID: S01-013
% Narrative Role: Optimization Deep-Dive
% Cognitive Objective: Understand why small integers share identical memory pointers.
% Plan Source: slide_plan.md / S01-013
% -----------------------------------------------------------------------------
\begin{frame}[fragile]{The CPython Integer Cache Optimization: \texttt{[-5, 256]}}
  \begin{columns}[T,onlytextwidth]
    \begin{column}{0.48\textwidth}
      \begin{CodeBlock}{Python}
# Within singleton cache [-5, 256]:
a = 256
b = 256
print(a is b)  # True (Shared singleton!)

# Outside cache (> 256):
x = 257
y = 257
print(x is y)  # False (Separate heap objects)
      \end{CodeBlock}
    \end{column}
    \begin{column}{0.48\textwidth}
      \begin{ConceptBox}{CPython Small Integer Singleton Pool}
        \scriptsize
        To avoid allocating millions of tiny objects, CPython pre-allocates an array of \InlineCode{PyLongObject} singletons for integers in the range $[-5, 256]$ at interpreter boot.
      \end{ConceptBox}
      \vspace{0.10cm}
      \begin{WarningBox}{Never Use 'is' for Numbers}
        \scriptsize Code relying on \InlineCode{x is 256} works by accident of CPython's cache and fails for \InlineCode{x is 257}. Always use \InlineCode{==}.
      \end{WarningBox}
    \end{column}
  \end{columns}
\end{frame}

% -----------------------------------------------------------------------------
% Slide ID: S01-014
% Narrative Role: Prediction / Comparison
% Cognitive Objective: Master the critical difference between mutating a heap object in-place and rebinding.
% Plan Source: slide_plan.md / S01-014
% -----------------------------------------------------------------------------
\begin{frame}[fragile]{In-Place Mutation (\texttt{+=}) vs Rebinding (\texttt{= +})}
  \begin{columns}[T,onlytextwidth]
    \begin{column}{0.48\textwidth}
      {\small\bfseries\color{AWSBlue}In-Place Mutation (\InlineCode{\_\_iadd\_\_})}\par\vspace{0.08cm}
      \begin{CodeBlock}{Python}
x = [1, 2]
y = x
x += [3]  # Calls __iadd__()
print("y:", y)  # [1, 2, 3] !
      \end{CodeBlock}
      \vspace{0.05cm}
      {\scriptsize\color{UpGradRedDark}Mutates the existing heap container in-place; \InlineCode{y} sees the modification.}
    \end{column}
    \hfill{\color{UpGradRule}\vrule width .5pt}\hfill
    \begin{column}{0.48\textwidth}
      {\small\bfseries\color{UpGradRed}Reference Rebinding (\InlineCode{\_\_add\_\_})}\par\vspace{0.08cm}
      \begin{CodeBlock}{Python}
a = [1, 2]
b = a
a = a + [3]  # Calls __add__()
print("b:", b)  # [1, 2] !
      \end{CodeBlock}
      \vspace{0.05cm}
      {\scriptsize\color{AWSEmerald}Allocates a fresh list and rebinds \InlineCode{a}; \InlineCode{b} remains untouched.}
    \end{column}
  \end{columns}
  \vspace{0.18cm}
  \TakeawayBanner{\InlineCode{+=} mutates existing heap RAM in-place; \InlineCode{= +} constructs a brand-new object.}
\end{frame}

% -----------------------------------------------------------------------------
% Slide ID: S01-015
% Narrative Role: Mechanism Reveal
% Cognitive Objective: Discover that default arguments are evaluated once when the module loads.
% Plan Source: slide_plan.md / S01-015
% -----------------------------------------------------------------------------
\begin{frame}[fragile]{The Root Cause: Function Definition-Time Evaluation}
  \begin{columns}[T,onlytextwidth]
    \begin{column}{0.48\textwidth}
      \begin{HighlightedCodeBlock}{Python}{1,7}
def append_event(event_id, log=[]):
    log.append(event_id)
    return log

# Inspecting the function's internal metadata:
print(append_event.__defaults__)
# Output: ([],) -> Persists across calls!
      \end{HighlightedCodeBlock}
    \end{column}
    \begin{column}{0.48\textwidth}
      \begin{ConceptBox}{Function Definition Lifecycle}
        \scriptsize
        \begin{enumerate}\setlength{\itemsep}{0.06cm}
          \item \textbf{Def Time (Module Load):} CPython compiles bytecode, allocates \InlineCode{[]} on heap, and stores pointer in \InlineCode{fn.\_\_defaults\_\_}.
          \item \textbf{Runtime Invocations:} Every call missing the argument reads from \InlineCode{fn.\_\_defaults\_\_[0]}!
        \end{enumerate}
      \end{ConceptBox}
    \end{column}
  \end{columns}
  \vspace{0.08cm}
  \TakeawayBanner{Default arguments are evaluated ONCE at function definition time, NOT at execution time.}
\end{frame}

% -----------------------------------------------------------------------------
% Slide ID: S01-016
% Narrative Role: Best Practice / Code Refactoring
% Cognitive Objective: Implement the standard Python idiom for safe default arguments.
% Plan Source: slide_plan.md / S01-016
% -----------------------------------------------------------------------------
\begin{frame}[fragile]{Production Fix: The Immutable Sentinel Pattern}
  \begin{columns}[T,onlytextwidth]
    \begin{column}{0.48\textwidth}
      \begin{tcolorbox}[enhanced,arc=1.2mm,boxrule=.4pt,colback=UpGradRed!8,colframe=UpGradRed,title={\bfseries\scriptsize Anti-Pattern (State Leak Hazard)},coltitle=white]
        \begin{CodeBlock}{Python}
def append_event(event_id, log=[]):
    log.append(event_id)
    return log
        \end{CodeBlock}
        \vspace{0.02cm}
        {\tiny\color{UpGradRedDark}*Shared mutable list causes cross-tenant state leaks.}
      \end{tcolorbox}
    \end{column}
    \begin{column}{0.48\textwidth}
      \begin{tcolorbox}[enhanced,arc=1.2mm,boxrule=.4pt,colback=AWSEmerald!8,colframe=AWSEmerald,title={\bfseries\scriptsize Production Pattern (Defensive Sentinel)},coltitle=white]
        \begin{CodeBlock}{Python}
def append_event(event_id, log=None):
    if log is None:
        log = [] # Fresh allocation per call!
    log.append(event_id)
    return log
        \end{CodeBlock}
        \vspace{0.02cm}
        {\tiny\color{AWSEmerald}*Allocates a fresh heap list on every call if log is None.}
      \end{tcolorbox}
    \end{column}
  \end{columns}
  \vspace{0.18cm}
  \TakeawayBanner{Always use immutable sentinels (\InlineCode{None}) for default collection parameters.}
\end{frame}

% -----------------------------------------------------------------------------
% Slide ID: S01-017
% Narrative Role: Concept / Memory Diagram
% Cognitive Objective: Visualize aliasing in nested compound structures.
% Plan Source: slide_plan.md / S01-017
% -----------------------------------------------------------------------------
\begin{frame}[fragile]{Object Aliasing: Two Names, One Compound Object}
  \begin{columns}[T,onlytextwidth]
    \begin{column}{0.38\textwidth}
      \begin{HighlightedCodeBlock}{Python}{2}
a = [1, [2, 3]]
b = a  # Pointer Assignment (Alias)

b.append(4)
print("a:", a) # [1, [2, 3], 4]!
      \end{HighlightedCodeBlock}
      \vspace{0.08cm}
      {\scriptsize\color{UpGradSlate}Assignment copies the outer pointer \InlineCode{0x100}. Zero cloning occurs.}
    \end{column}
    \begin{column}{0.58\textwidth}
      \centering
      \DrawCopySemanticsDiagram{1}
    \end{column}
  \end{columns}
  \vspace{0.15cm}
  \TakeawayBanner{Variable assignment (\InlineCode{b = a}) creates an alias to the identical outer container.}
\end{frame}

% -----------------------------------------------------------------------------
% Slide ID: S01-018
% Narrative Role: Mechanistic Deep-Dive
% Cognitive Objective: Understand that shallow copy duplicates ONLY the top-level pointer array.
% Plan Source: slide_plan.md / S01-018
% -----------------------------------------------------------------------------
\begin{frame}[fragile]{Shallow Copy: What Actually Gets Duplicated?}
  \begin{columns}[T,onlytextwidth]
    \begin{column}{0.38\textwidth}
      \begin{HighlightedCodeBlock}{Python}{2,4}
a = [1, [2, 3]]
b = a.copy()  # Shallow Copy

b[1].append(99) # Mutate inner list!
print("a:", a)  # [1, [2, 3, 99]]!
      \end{HighlightedCodeBlock}
      \vspace{0.08cm}
      {\scriptsize\color{AWSOrange}\textbf{Warning:} The outer container is duplicated, but inner pointers are shared!}
    \end{column}
    \begin{column}{0.58\textwidth}
      \centering
      \DrawCopySemanticsDiagram{2}
    \end{column}
  \end{columns}
  \vspace{0.15cm}
  \TakeawayBanner{Shallow copy duplicates the outer array, but retains shared pointers to nested objects.}
\end{frame}

% -----------------------------------------------------------------------------
% Slide ID: S01-019
% Narrative Role: Solution / Diagram
% Cognitive Objective: Master recursive graph duplication via copy.deepcopy.
% Plan Source: slide_plan.md / S01-019
% -----------------------------------------------------------------------------
\begin{frame}[fragile]{Deep Copy: Recursive Pointer Graph Duplication}
  \begin{columns}[T,onlytextwidth]
    \begin{column}{0.38\textwidth}
      \begin{HighlightedCodeBlock}{Python}{2}
import copy
b = copy.deepcopy(a)  # Deep Clone

b[1].append(99)
print("a:", a) # [1, [2, 3]] (Intact!)
print("b:", b) # [1, [2, 3, 99]]
      \end{HighlightedCodeBlock}
      \vspace{0.08cm}
      {\scriptsize\color{AWSEmerald}\textbf{Complete Isolation:} Recursively duplicates all nested compound objects.}
    \end{column}
    \begin{column}{0.58\textwidth}
      \centering
      \DrawCopySemanticsDiagram{3}
    \end{column}
  \end{columns}
  \vspace{0.15cm}
  \TakeawayBanner{Use \InlineCode{copy.deepcopy()} to ensure complete decoupling of nested compound structures.}
\end{frame}

% -----------------------------------------------------------------------------
% Slide ID: S01-020
% Narrative Role: VM Internals
% Cognitive Objective: Understand immediate deallocation when ob_refcnt hits zero.
% Plan Source: slide_plan.md / S01-020
% -----------------------------------------------------------------------------
\begin{frame}[fragile]{Memory Deallocation: Reference Counting Mechanics}
  \begin{columns}[T,onlytextwidth]
    \begin{column}{0.48\textwidth}
      \begin{HighlightedCodeBlock}{Python}{2,5,8}
import sys
x = [10, 20]
# refcnt = 1 (+1 temp arg = 2)
print(sys.getrefcount(x)) # 2

y = x
# refcnt = 2 (+1 temp arg = 3)
print(sys.getrefcount(x)) # 3

del x; del y
# ob_refcnt -> 0 ==> Immediate free()!
      \end{HighlightedCodeBlock}
    \end{column}
    \begin{column}{0.48\textwidth}
      \begin{ConceptBox}{Deterministic Deallocation}
        \scriptsize
        Python's primary memory management engine is \textbf{deterministic reference counting}:
        \begin{itemize}\setlength{\itemsep}{0.08cm}
          \item Every pointer assignment increments \InlineCode{ob\_refcnt}.
          \item Rebinding or \InlineCode{del} decrements \InlineCode{ob\_refcnt}.
          \item When \InlineCode{ob\_refcnt == 0}, CPython immediately calls C \InlineCode{free()}.
        \end{itemize}
      \end{ConceptBox}
    \end{column}
  \end{columns}
  \vspace{0.15cm}
  \TakeawayBanner{CPython frees heap memory the exact moment an object's reference count reaches zero.}
\end{frame}

% -----------------------------------------------------------------------------
% Slide ID: S01-021
% Narrative Role: Failure Mode / Edge Case
% Cognitive Objective: Understand why reference counting alone fails when objects point to each other.
% Plan Source: slide_plan.md / S01-021
% -----------------------------------------------------------------------------
\begin{frame}[fragile]{The Failure Mode: Cyclic Reference Traps}
  \begin{columns}[T,onlytextwidth]
    \begin{column}{0.48\textwidth}
      \begin{HighlightedCodeBlock}{Python}{2,3}
lst = []
lst.append(lst) # Self-referential cycle!
del lst         # Stack pointer removed...
      \end{HighlightedCodeBlock}
      \vspace{0.10cm}
      \begin{WarningBox}{Reference Count Never Reaches 0}
        \scriptsize
        Even though \InlineCode{lst} is deleted from the stack, the heap object references itself (\InlineCode{ob\_refcnt = 1}). It is unreachable from user code, yet cannot be freed by refcounting!
      \end{WarningBox}
    \end{column}
    \begin{column}{0.48\textwidth}
      \begin{center}
        \begin{tikzpicture}[font=\scriptsize\bfseries]
          \node[draw=UpGradRed,fill=UpGradMist,rounded corners=2mm,thick,minimum width=2.6cm,minimum height=1.2cm,align=center] (nodeA) at (0,0) {
            Heap Node A\\ \tiny ob\_refcnt = 1
          };
          \draw[upgrad/pointer,UpGradRed,thick] (nodeA.north east) to[out=45,in=135,distance=1.6cm] (nodeA.north west);
          \node[font=\tiny\color{UpGradRedDark}] at (0,1.2) {Self-Reference};
        \end{tikzpicture}
      \end{center}
      \begin{ConceptBox}{Generational Cyclic GC}
        \scriptsize CPython runs a secondary periodic cyclic garbage collector that tracks container pointers and breaks isolated reference cycles.
      \end{ConceptBox}
    \end{column}
  \end{columns}
\end{frame}

% -----------------------------------------------------------------------------
% Slide ID: S01-022
% Narrative Role: Applied Production Code
% Cognitive Objective: Apply pointer semantics and sentinels to build an enterprise-grade batch pipeline.
% Plan Source: slide_plan.md / S01-022
% -----------------------------------------------------------------------------
\begin{frame}[fragile]{Authentic Application: Multi-Tenant Batch Ingestion}
  \begin{HighlightedCodeBlock}{Python}{4,8}
import copy

class BatchIngestor:
    def __init__(self, tenant_id: str, default_meta: dict | None = None):
        self.tenant_id = tenant_id
        # Defensively clone mutable configs to prevent cross-tenant state corruption!
        self.meta = copy.deepcopy(default_meta) if default_meta is not None else {}
        self.records: list[dict] = []

    def ingest(self, batch: list[dict]) -> None:
        # Shallow copy batch container to protect against external mutation during write
        self.records.extend(batch.copy())
  \end{HighlightedCodeBlock}
  \vfill
  \TakeawayBanner{Defensive deep copying on config and shallow copying on batch inputs ensures memory isolation.}
\end{frame}

% -----------------------------------------------------------------------------
% Slide ID: S01-023
% Narrative Role: Interactive Exercise
% Cognitive Objective: Test student mastery of shallow copy mechanics on nested dictionaries.
% Plan Source: slide_plan.md / S01-023
% -----------------------------------------------------------------------------
\begin{frame}[fragile]{Transfer Challenge 1: The Nested Dictionary Mutation Riddle}
  \begin{columns}[T,onlytextwidth]
    \begin{column}{0.54\textwidth}
      \begin{CodeBlock}{Python}
d1 = {'meta': [10, 20], 'id': 1}
d2 = d1.copy()  # Shallow Copy

d2['meta'].append(30)
d2['id'] = 2

print("d1['meta']:", d1['meta'])
print("d1['id']:", d1['id'])
      \end{CodeBlock}
    \end{column}
    \begin{column}{0.43\textwidth}
      \begin{QuestionBox}{Predict the Output for d1}
        \begin{itemize}\setlength{\itemsep}{0.12cm}
          \item[\textbf{A)}] \InlineCode{'meta': [10, 20], 'id': 1}
          \item[\textbf{B)}] \InlineCode{'meta': [10, 20, 30], 'id': 1}
          \item[\textbf{C)}] \InlineCode{'meta': [10, 20, 30], 'id': 2}
          \item[\textbf{D)}] \InlineCode{'meta': [10, 20], 'id': 2}
        \end{itemize}
      \end{QuestionBox}
      \vspace{0.10cm}
      {\scriptsize\color{UpGradSlate}Why did one key change while the other remained untouched?}
    \end{column}
  \end{columns}
\end{frame}

% -----------------------------------------------------------------------------
% Slide ID: S01-024
% Narrative Role: Solution / Explanation
% Cognitive Objective: Reinforce that primitive dictionary values are rebound while mutable nested lists are mutated in-place.
% Plan Source: slide_plan.md / S01-024
% -----------------------------------------------------------------------------
\begin{frame}[fragile]{Transfer Challenge 1: Solution \& Pointer Trace}
  \begin{columns}[T,onlytextwidth]
    \begin{column}{0.48\textwidth}
      \begin{TerminalWindow}[PYTHON EVALUATION]
>>> d1['meta']: [10, 20, 30]  \# Mutated in-place! \\
>>> d1['id']: 1               \# Unchanged!
      \end{TerminalWindow}
      \vspace{0.12cm}
      \begin{ConceptBox}{Option B is Correct!}
        \scriptsize
        \begin{itemize}
          \item \InlineCode{d2['id'] = 2} repoints \InlineCode{d2}'s key to integer \InlineCode{2}.
          \item \InlineCode{d2['meta']} resolves the shared heap list pointer and appends \InlineCode{30} in-place!
        \end{itemize}
      \end{ConceptBox}
    \end{column}
    \begin{column}{0.48\textwidth}
      \begin{center}
        \begin{tikzpicture}[font=\ttfamily\tiny,scale=0.88,transform shape]
          \node[draw=UpGradNight,fill=white,rounded corners=1mm] (d1) at (0,2.0) {d1: {'meta': 0x50, 'id': 1}};
          \node[draw=AWSBlue,fill=white,rounded corners=1mm] (d2) at (0,0.6) {d2: {'meta': 0x50, 'id': 2}};
          \node[draw=AWSEmerald,fill=AWSEmerald!15,rounded corners=1mm,thick] (lst) at (3.6,1.3) {0x50: [10, 20, 30]};
          \draw[upgrad/pointer,UpGradRed] (d1.east) -- (lst.west);
          \draw[upgrad/pointer,AWSBlue] (d2.east) -- (lst.west);
        \end{tikzpicture}
      \end{center}
    \end{column}
  \end{columns}
  \vspace{0.06cm}
  \TakeawayBanner{Primitives inside shallow copies rebind safely; nested mutable collections alias.}
\end{frame}

% -----------------------------------------------------------------------------
% Slide ID: S01-025
% Narrative Role: Synthesis
% Cognitive Objective: Consolidate the 4 core rules of Python memory.
% Plan Source: slide_plan.md / S01-025
% -----------------------------------------------------------------------------
\begin{frame}{Block 1 Synthesis: The Pointer Mental Model}
  \begin{columns}[T,onlytextwidth]
    \begin{column}{0.48\textwidth}
      \begin{ConceptBox}{1. Variables are Pointers}
        \scriptsize Stack names store 64-bit heap addresses. Variables do not hold raw data.
      \end{ConceptBox}
      \vspace{0.08cm}
      \begin{ConceptBox}{2. Defaults Evaluated Once}
        \scriptsize Stored in \InlineCode{fn.\_\_defaults\_\_} at module load. Always use \InlineCode{None} sentinels.
      \end{ConceptBox}
    \end{column}
    \begin{column}{0.48\textwidth}
      \begin{ConceptBox}{3. Rebinding vs Mutation}
        \scriptsize \InlineCode{=} moves the pointer; \InlineCode{.append()} mutates the heap object in-place.
      \end{ConceptBox}
      \vspace{0.08cm}
      \begin{ConceptBox}{4. Copy Hierarchy}
        \scriptsize Assignment = alias; Shallow = outer only; Deep = recursive graph clone.
      \end{ConceptBox}
    \end{column}
  \end{columns}
  \vspace{0.10cm}
  \TakeawayBanner{Mastering Python memory begins with tracking what pointers reference on the heap.}
\end{frame}

% -----------------------------------------------------------------------------
% Slide ID: S01-026
% Narrative Role: Transition / Bridge
% Cognitive Objective: Transition from individual memory pointers to collections of millions of pointers.
% Plan Source: slide_plan.md / S01-026
% -----------------------------------------------------------------------------
\begin{frame}{Bridge to Block 2: The Computational Cost of Containers}
  \begin{center}
    \vspace{0.15cm}
    \begin{tcolorbox}[enhanced,arc=2mm,boxrule=.6pt,colback=UpGradNight,colframe=AWSSquid,
      left=6mm,right=6mm,top=4mm,bottom=4mm,width=0.92\linewidth]
      {\UpGradDisplayFont\fontsize{12}{16}\selectfont\bfseries\color{white}
        “If every variable is merely a pointer to a heap object, why do different collection containers take vastly different amounts of time to find the exact same pointer?”
      }
    \end{tcolorbox}
  \end{center}
  \vfill
  \begin{columns}[c,onlytextwidth]
    \begin{column}{0.48\textwidth}
      \StatCard{UpGradRed}{120 ms}{List Linear Scan}{Scanning 10,000,000 pointers one-by-one}
    \end{column}
    \begin{column}{0.48\textwidth}
      \StatCard{AWSEmerald}{45 ns}{Hash Table Lookup}{Direct address jump via hash() \% N}
    \end{column}
  \end{columns}
\end{frame}
